\section{傅里叶级数的渐进特性，吉布斯现象}\label{sec:Asymptotic_Behaviour}

\subsection{收敛速度的估计}\label{sub:converge_speed}
在用计算机模拟函数的傅里叶级数展开时，只能取有限项，自然要问计算到多少项时误差足够小，为此，我们不加证明地给出以下定理\cite{brad_osgood_ee261}：
\begin{theorem}
    设$f\in C^p(\mathbb{R} )(p\geq1)$是周期函数，其傅里叶系数为$c_k$，部分和
    \[S_N^f(t)=\sum_{k=-N}^{N}c_k \mathrm{e}^{ik\omega t}\]
    在$\mathbb{R} $上逐点收敛、内闭一致收敛，且
    \[\max\left|f(t)-S_N^f(t)\right|<\frac{1}{N^{p-\frac{1}{2}}}\]
其中$C^p(\mathbb{R} )$表示p次连续可导的函数集。
\end{theorem}

\subsection{吉布斯现象}\label{sub:gibbs_phenomenon}
当$f(t)$不连续时，傅里叶级数的会在间断点处产生\textbf{吉布斯现象} (Gibbs' Phenomenon)：
\begin{definition}[吉布斯现象]
    部分和$S_N^f(t)$在间断点处总会\textbf{过冲}（在间断点两侧出现超过原函数的峰值），过冲幅度约为跳变幅度的9\%，并且$S_N^f(t)$会在间断点附近高频振荡
\end{definition}
\begin{example}
    对于跳变幅度为2、周期为$2\pi$的周期矩形脉冲信号
\[R(x) =
    \begin{cases}
        1  & \text{if } 0<x<\pi  \\
        -1 & \text{if } -\pi<x<0
    \end{cases}\]
其傅里叶级数的跳变值为1，$\varlimsup_{N \to \infty}S_N^R(t)=1.089490 \dots$。这是因为光滑的基函数很难逼近这种剧烈的局部变化，不得不用高频分量来补偿，高频分量带来了剧烈震动。
\end{example}
$\varlimsup_{N \to \infty}S_N^R(t)>1$并不意味着狄利克雷定理失效，因为定理给出的是逐点收敛：
\[\varlimsup_{N \to \infty}S_N^R(t)=\lim_{N \to \infty}\sup_{t\in \mathbb{R} }S_N^R(t)\neq \sup_{t\in\mathbb{R}}\lim_{N\to\infty}S_N^R(t)\]
使极限号与取上界号交换的一个充分条件是一致收敛，吉布斯现象表明，傅里叶级数在包含间断点的区间上不可能一致收敛。
\begin{figure}[H]
    \centering
    \includegraphics[width=0.8\textwidth]{gibbs}
    \caption{吉布斯现象示意图}
\end{figure}

为了直观地理解它，我们来看一个经典的例子：
\begin{example}
    \begin{align*}
    f_n(x)=
    \begin{cases}
        nx   & \text{if } 0<x\leq 1/n        \\
        2-nx & \text{if } 1/n<x<2/n \\
        0    & \text{otherwise}
    \end{cases}
\end{align*}
随n增大，$f(x)$逐点趋于0，因为对每一点$2/n$总能取到更小的值；但$f(x)$的最大值永远是1。
\end{example}

\subsection{狄利克雷核}\label{sub:dirichlet_kernel}
本小节引入狄利克雷核并介绍它的一些性质，并在下一小节用它来辅助我们对傅里叶级数收敛性的讨论。直接由狄利克雷条件证明收敛性是很困难的，并且需要更加专业的分析学工具，我们将给出更强的条件下的证明，并从这个更强的条件获得一些额外的性质。

研究傅里叶级数的渐进特性时，一个非常好用的工具是\textbf{狄利克雷核} (Dirichlet kernel)：
\begin{definition}[狄利克雷核]
    \[D_N(t)=\sum_{k=-N}^{N}\mathrm{e}^{\mathrm{i}k\omega t}=1+\sum_{k=1}^{N}\left(\mathrm{e}^{\mathrm{i}k\omega t}+\mathrm{e}^{-\mathrm{i}k\omega t}\right)=1+2\sum_{k=1}^{N}\cos(k\omega t)\]
    它是依赖于所研究函数的周期T的，但简便起见，在符号$D_N(t)$中不体现这一点。
\end{definition}

我们可以用等比数列求和或积化和差裂项的方法化简$D_N(t)$：
\begin{align*}
    D_N(t) & =\sum_{k=-N}^{N}\mathrm{e}^{\mathrm{i}k\omega t}=\mathrm{e}^{-\mathrm{i}N\omega t}\frac{1-\mathrm{e}^{\mathrm{i}(2N+1)\omega t}}{1-\mathrm{e}^{\mathrm{i}\omega t}}                  \\
           & =\frac{\mathrm{e}^{\mathrm{i}(N+1)\omega t}-\mathrm{e}^{-\mathrm{i}N\omega t}}{\mathrm{e}^{\mathrm{i}\omega t}-1}                                   \\
    D_N(t) & =1+\sum_{k=1}^{N}\left(\mathrm{e}^{\mathrm{i}k\omega t}+\mathrm{e}^{-\mathrm{i}k\omega t}\right)=1+2\sum_{k=1}^{N}\cos(k\omega t)           \\
           & =1+\frac{2}{\sin(\frac{\omega t}{2})}\sum_{k=1}^{N}\cos(k\omega t)\sin\left(\frac{\omega t}{2}\right)                   \\
           & =1+\frac{1}{\sin(\frac{\omega t}{2})}\sum_{k=1}^{N}\left[\sin\left(\left(k+\frac{1}{2}\right)\omega t\right)-\sin\left(\left(k-\frac{1}{2}\right)\omega t\right)\right] \\
           & =1+\frac{\sin\left((N+\frac{1}{2})\omega t\right)-\sin(\frac{\omega t}{2})}{\sin(\frac{\omega t}{2})}                     \\
           & =\frac{\sin\left((N+\frac{1}{2})\omega t\right)}{\sin(\frac{\omega t}{2})}
\end{align*}
这两种结果是相符的，读者可自行验证，并且可以从后一结果想象出狄利克雷核的函数图像，它被$\pm 1/\sin(\frac{\omega t}{2})$包络并高速振荡。函数图像如下。
\begin{figure}[H]
    \centering
    \includegraphics[width=0.8\textwidth]{Figure_3}
    \caption{狄利克雷核$D_N(t)$的函数图像}
\end{figure}

傅里叶级数的部分和是狄利克雷核与被展开函数的卷积（周期函数的卷积，见\ref{sub:conv_of_period_func}）：
\begin{proposition}
    \[S_N^f(t)=\frac{1}{T}\int_{T}f(\tau)D_N(t-\tau)\,d\tau=\frac{1}{T}f*D_N(t)\]
\end{proposition}
\begin{proof}
    \begin{align*}
    S_N^f(t) & =\sum_{-N}^{N}c_k \mathrm{e}^{\mathrm{i}k\omega t}\\
             & =\sum_{-N}^{N}\left(\frac{1}{T}\int_{T}f(\tau)\mathrm{e}^{-\mathrm{i}k\omega \tau}\,d\tau\right) \mathrm{e}^{\mathrm{i}k\omega t}\\
             & =\frac{1}{T}\int_{T}f(\tau)\sum_{k=-N}^{N}\mathrm{e}^{\mathrm{i}k\omega (t-\tau)}\,d\tau\\
             & =\frac{1}{T}\int_{T}f(\tau)D_N(t-\tau)\,d\tau\\
             & =\frac{1}{T}\int_{T}f(t-\tau)D_N(\tau)\,d\tau                                                & (\tau\to t-\tau)\\
             & =\frac{1}{T}\int_{T}f(t+\tau)D_N(\tau)\,d\tau                                                & (\tau\to t+\tau)
\end{align*}
\end{proof}

在讨论傅里叶级数的收敛性前，先给出两个引理。第一个引理表明狄利克雷核在半周期上
积分值为$\frac{T}{2}$，在证明傅里叶级数的逐点收敛性时将用到它。
\begin{lemma}
    $$\int_{-\frac{T}{2}}^{0}D_N(t)\,dt=\int_{0}^{\frac{T}{2}}D_N(t)\,dt=\frac{T}{2}$$
\end{lemma}
\begin{proof}
    \begin{align*}
    D_N(t)&=1+2\sum_{k=1}^{N}\cos(k\omega t)\\
    \int_{0}^{\frac{T}{2}}D_N(t)\,dt & =\int_{0}^{\frac{T}{2}}\left(1+\sum_{k=1}^{N}\cos(k\omega t)\right)\,dt                              \\
                                     & =\frac{T}{2}+\sum_{k=1}^{N}\left.\frac{\sin(k\omega t)}{k\omega}\right|_{0}^{\frac{T}{2}} \\
                                     & =\frac{T}{2}+\frac{1}{\omega}\sum_{k=1}^{N}\frac{\sin(k\pi)}{k}=\frac{T}{2}               
\end{align*}
$D_N(t)$是偶函数，得证。
\end{proof}

\begin{lemma}[贝塞尔不等式]
设$ f\in L^2([0,T]),c_n=\dfrac{1}{T}\int_{T}f(t)\mathrm{e}^{-\mathrm{i}k\omega t}\,dt$，则
\begin{align*}
    \sum_{-\infty}^{\infty}|c_n|^2\leq \frac{1}{T}\int_{T}|f(t)|^2\,dt=\frac{1}{T}\|f\|_2^2
\end{align*}
\end{lemma}
它给出了傅里叶系数平方和的上界的估计。由此可以看出$c_n\to 0,n\to\infty$.
\begin{proof}
    \begin{align*}
    \left|f(t)-\sum_{n=-N}^{N}c_n \mathrm{e}^{\mathrm{i}n\omega t}\right|^2 & =\left(f(t)-\sum_{n=-N}^{N}c_n \mathrm{e}^{\mathrm{i}n\omega t}\right)\left(f(t)-\sum_{n=-N}^{N}c_n \mathrm{e}^{\mathrm{i}n\omega t}\right)^*\\
    & =\left(f(t)-\sum_{n=-N}^{N}c_n \mathrm{e}^{\mathrm{i}n\omega t}\right)\left(f^*(t)-\sum_{n=-N}^{N}c_n \mathrm{e}^{-\mathrm{i}n\omega t}\right)\\
    & =|f(t)|^2-\sum_{n=-N}^{N}\left(c_n^*f(t)\mathrm{e}^{\mathrm{i}n\omega t}+c_n f^*(t)\mathrm{e}^{-\mathrm{i}n\omega t}\right)+\sum_{m,n=-N}^{N}c_m c_n^*\mathrm{e}^{\mathrm{i}(m-n)\omega t}
\end{align*}
将上式在一个周期上积分，我们知道
\[\int_{T}f(t)\mathrm{e}^{\mathrm{i}n\omega t}\,dt=Tc_n,\int_{T}\mathrm{e}^{\mathrm{i}(m-n)\omega t}\,dt=\begin{cases}
        0 & \text{if }m\neq n \\
        T & \text{if }m=n
    \end{cases}\]
故\begin{align*}
      & \int_{T}|f(t)|^2\,dt-\sum_{n=-N}^{N}\left(c_n^*\int_{T}f(t)\mathrm{e}^{\mathrm{i}n\omega t}\,dt+c_n \int_{T}f^*(t)\mathrm{e}^{-\mathrm{i}n\omega t}\,dt\right)+\sum_{m,n=-N}^{N}c_m c_n^*\int_{T}\mathrm{e}^{\mathrm{i}(m-n)\omega t}\,dt \\
    = & \int_{T}|f(t)|^2\,dt-T\sum_{n=-N}^{N}(c_n^* c_n+c_n c_n^*)+T\sum_{n=-N}^{N}c_n^* c_n\\
    = & \int_{T}|f(t)|^2\,dt-T\sum_{n=-N}^{N}|c_n|^2
\end{align*}
这是非负函数的积分，积分值非负，即\begin{align*}
    \sum_{-\infty}^{\infty}|c_n|^2\leq \frac{1}{T}\int_{T}|f(t)|^2\,dt<\infty
\end{align*}
\end{proof}

\subsection{分段光滑条件下的一致收敛性}\label{sub:uniform_converge}
直接由狄利克雷条件证明逐点收敛性需要很专业的分析学工具，但我们可以适当地加强狄利克雷条件，以得到更直观、简单的结论\cite{folland}。回忆\ref{sub:fourier_series_convergency}中对于分段光滑的定义：
\begin{definition}[分段光滑]
在区间$[a,b]$上，函数$f$称为\textbf{分段光滑} (piecewise smooth)，如果存在区间的划分$a=x_0<x_1<\dots <x_n=b$，使得在每个子区间$(x_{i-1},x_i)$上f可微，且区间端点处$f$的左右导数存在，记作$f\in PS[a,b]$。
\end{definition}
我们研究的多数函数是满足这样的性质的，并且我们将看到满足此条件会带来一些额外的性质。此时就可以相对简单地证明逐点收敛性：
\begin{theorem}[傅里叶级数的逐点收敛性]
    \[\lim_{N\to\infty}S_N^f(t_0)=\frac{f(t_0+)+f(t_0-)}{2}\]
\end{theorem}
\begin{proof}
\begin{align*}
    S_N^f(t_0)-\frac{f(t_0^+)+f(t_0^-)}{2} =&\frac{1}{T}\left(\int_{T}f(t_0-\tau)D_N(\tau)\,d\tau\right.\\
    &\left.-\int_{0}^{\frac{T}{2}}f(t_0^+)D_N(\tau)\,d\tau-\int_{-\frac{T}{2}}^{0}f(t_0^-)D_N(\tau)\,d\tau\right)\\
    =&\frac{1}{T}\left(\int_{0}^{\frac{T}{2}}(f(t_0-\tau)-f(t_0^+))D_N(\tau)\,d\tau\right.\\
    &+\left.\int_{-\frac{T}{2}}^{0}(f(t_0-\tau)-f(t_0^-))D_N(\tau)\,d\tau\right)\\
    S_N^f(t_0)-\frac{f(t_0^+)+f(t_0^-)}{2}=&\frac{1}{T}\int_{T}g(t)\left(\mathrm{e}^{\mathrm{i}(N+1)\omega t}-\mathrm{e}^{\mathrm{i}N\omega t}\right)\,dt                            
\end{align*}
其中\begin{align*}
g(t):=\begin{cases}
\dfrac{f(t_0+t)-f(t_0^-)}{\mathrm{e}^{\mathrm{i}\omega t}-1} & \text{if }-\frac{T}{2}<t_0<0 \\
\dfrac{f(t_0+t)-f(t_0^+)}{\mathrm{e}^{\mathrm{i}\omega t}-1} & \text{if }0<t_0<\frac{T}{2}
\end{cases}
\end{align*}                        
由洛必达法则，$t\to 0$时，\[\lim_{t\to 0^+}g(t)=\lim_{t\to 0^+}\frac{f(t_0+t)-f(t_0^+)}{\mathrm{e}^{\mathrm{i}\omega t}-1}=\lim_{t\to 0^+}\frac{f'(t_0+t)}{\mathrm{i}\omega \mathrm{e}^{\mathrm{i}\omega t}}=\lim_{t\to 0^+}\frac{f'(t_0^+)}{\mathrm{i}\omega}\]
$t\to 0^-$时同理。故g分段连续，当然是平方可积的，由贝塞尔不等式，$g(t)$的傅里叶系数平方和收敛，通项趋于0，$S_N^f(t_0)-\left(f(t_0^+)+f(t_0^-)\right)/2=C_{-(N+1)}-C_N\to 0$，得证。
\end{proof}

在分段光滑的条件下，容易得到$f'(t)$的傅里叶系数，注意微积分基本定理可以分区间使用：
\begin{proposition}[导函数的傅里叶系数]
    \begin{equation*}
    a'_n=n\omega b_n,b'_n=-n\omega a_n,c'_n=\mathrm{i}n\omega c_n
\end{equation*}
\end{proposition}
\begin{proof}
    以$c_n$为例：
\begin{align*}
    c'_n & =\frac{1}{T}\int_{T}f'(t)\mathrm{e}^{-\mathrm{i}n\omega t}\,dt\\
         & =\frac{1}{T}\evalat{f(t)\mathrm{e}^{-\mathrm{i}n\omega t}}{0}{T}+\mathrm{i}n\omega\int_{T}f(t)\mathrm{e}^{-\mathrm{i}n\omega t}\,dt=\mathrm{i}n\omega c_n
\end{align*}
\end{proof}

f的原函数F的傅里叶系数同理，并且只要\textbf{分段连续}（见\ref{sub:fourier_series_convergency}）即可保证f可积，但是我们必须保证F是周期函数，这要求f的直流分量为0：
\[F(t+T)-F(t)=\int_{T}f(t)dt=Tc_0=0,c_0=0\]
此时，用刚刚得到的公式（2.16）就可直接得到F的傅里叶系数：
\begin{proposition}[原函数的傅里叶系数]
    \begin{equation*}
    A_n=\frac{a_n}{n\omega},B_n=\frac{b_n}{n\omega},C_n=\frac{c_n}{\mathrm{i}n\omega}
\end{equation*}
\end{proposition}

分段光滑还能推出傅里叶级数的\textbf{一致收敛性}
\begin{theorem}
$f\in PS([0,T])\implies S_N^f(t)\rightrightarrows f(t)$
\end{theorem}
以下证明中，我们将使用魏尔斯特拉斯M判别法：
\begin{theorem}[M判别法]
    对于函数项级数$\sum_{n=1}^{\infty}f_n(x)$，如果存在正项级数$\sum_{n=1}^{\infty}M_n<\infty$使得在区间E上$|f_n(x)|<M_n$，则$\sum_{n=1}^{\infty}f_n(x)$在E上绝对收敛且一致收敛。
\end{theorem}
\begin{proof}
    要证明一致收敛性，只需证明
    $$\left|\sum_{n=-\infty}^{\infty}c_n \mathrm{e}^{\mathrm{i}n\omega t}\right|\leq \sum_{n=-\infty}^{\infty}|c_n|<\infty$$
    记$c'_n$为$f'$的傅里叶系数，$c'_n =\mathrm{i}n\omega c_n$，
    \begin{align*}
    \sum_{n=-\infty}^{\infty}|c_n|=|c_0|+\sum_{n\neq 0}\left| \frac{c'_n}{n}\right|\leq & |c_0|+\left(\sum_{n\neq 0}\frac{1}{n^2}\right)^{1/2}\left(\sum_{n\neq 0}|c'_n|^2\right)^{1/2}<\infty
\end{align*}
最后一步使用了柯西-施瓦兹不等式。
\end{proof}

请读者思考：我们探究了指数形式傅里叶级数收敛的条件，对于三角函数形式的傅里叶级数应该怎么办？
